\section{A falsifiable research program}

The proposed agenda has a narrow causal core and a wider commercial test. The causal core asks whether an aftermarket harness improves completion with the model held fixed. The commercial test asks whether that improvement remains valuable after expert nudging, deterministic software, implementation costs, and provider progress are counted. Confusing the two leads to a common error: a technically interesting harness can demonstrate recovered tasks without becoming the best use of a non-frontier research budget.

\textbf{H1: Productized intervention recovers permitted stops.} For a preregistered set of legitimate tasks that a stock product leaves incomplete, the same model inside a bounded specialist harness reaches more independently verified terminal states than the stock product under a specified budget. Reject H1 if the uplift disappears on held-out services or if the apparent gain comes from a weaker acceptance check. This hypothesis does not say every observed refusal is recoverable.

\textbf{H2: A harness can reduce the nudge tax.} Against a skilled human nudging the stock product, the harness reduces user-active tactical interventions while preserving necessary approvals and accepted quality. Reject H2 if the owner must still repair outputs, supply context repeatedly, or handle more exceptions than before. Record waiting time separately so reduced active minutes are not mistaken for a faster service.

\textbf{H3: The value lies in system design, not merely more calls.} In ablations with matched model and tool budgets, task specification, durable state, allowed connectors, and verification each contribute identifiable gains or reduce distinct failure classes. Reject the broad harness story if a short prompt template or ordinary deterministic connector explains nearly all uplift. The SWE-agent result makes interface design plausible as an experimental variable, but its coding-task setting does not settle this cross-domain hypothesis \citep{sweagent}.

\textbf{H4: The commercial wedge survives change.} After a provider releases a stronger native product, the specialist retains enough advantage in a chosen domain to cover onboarding, operation, maintenance, and sales. Reject H4 if native completion converges, if platform access disappears, or if the cost per accepted task remains above the customer's alternative. The relevant test is longitudinal. A one-time frontier gap is a launch window, not automatically a durable company.

A practical study sequence begins with twenty to fifty consented task templates in one domain, each with a clear owner and external verifier. The precise count is a proposed study design, not a claimed sample or power calculation; power analysis should follow pilot variance. First collect stock-product and human-nudging episodes. Next build the minimum harness that addresses the most common recoverable stop, then add components only when traces identify a bottleneck. Keep a holdout for each service and task family. Finally repeat the study after a provider update. The outcome is a map of which interventions matter under which conditions, not a generic agent benchmark score.

The superlative thesis needs a broader comparison set. Allocate the same engineering budget to at least two alternatives: improving a domain model or prompt-only product, and building conventional integration software without an agent loop. Define a shared customer outcome metric and compare value per engineering dollar, adjusted for human work and operating risk. The aftermarket position is rejected for that cohort if an alternative consistently produces more accepted work per dollar. It can still hold in another domain, but the claim must be narrowed rather than rescued by changing the denominator after seeing results.

Some evidence is already suggestive but not dispositive. Anthropic's engineering reports show that harness choices can change long-running coding behavior, and a later report explicitly removed a previously useful scaffold as the underlying model changed \citep{anthropiclong,anthropic2026}. Berkeley's compound-system argument and the existence of provider SDKs show why a single model score is an incomplete description of an agent product \citep{compound,openaiagents}. These sources also strengthen the counterargument that frontier providers may internalize the opportunity. The proposed studies are designed to distinguish a durable specialist advantage from a temporary implementation gap.

Research ethics and customer trust belong in the protocol but need not dominate the thesis. Tasks must be owner-authorized and compatible with the third-party route used. Properly prohibited requests are not in the recoverable-stop target. Human approvals remain explicit; the harness cannot silently turn an observed decline into an unbounded search for a more permissive path. Evaluation data should avoid public credentials and private customer traces. These boundaries make the commercial claim harder to win, and more meaningful if it does win.

\subsection*{What would turn a research result into a company?}

An uplift in accepted outcomes is necessary but insufficient for a commercial conclusion. A prospective buyer needs to see the same task family recur, the value of an accepted result exceed the total cost of delivery, and the system work under that buyer's actual authorization and support process. The research report should therefore publish a deployment envelope: task volume, tolerated delay, required human decisions, integration work, ordinary failure rate, and the conditions under which the workflow is withdrawn. A positive result in one envelope should not be advertised as a general agent capability.

The specialist also has to know who pays for the hidden work. If the vendor maintains service adapters, monitors drift, resolves uncertain remote writes, and staffs escalations, those are vendor costs even when a model API bill is small. If the customer supplies skilled reviewers or manually repairs records, those are customer costs even when the vendor dashboard shows success. The commercial test should price both sides and report the split. A bargain to one party may simply transfer the nudge tax to another.

Finally, the customer may value a narrower guarantee more than higher average completion. A workflow that finishes eight of ten tasks but cannot identify the other two may be harder to buy than one that finishes fewer tasks and reliably escalates the rest with evidence. A purchasing study should therefore ask buyers what terminal states they can operationally accept: completed, pending external confirmation, handed to a named owner, or stopped with a reason. The answer may differ by task family and consequence. This is where workflow research meets product design: a harness has to make noncompletion legible enough to be handled, rather than hiding it behind an optimistic summary.
